\section{Parabolic norms and uniform local estimates}\label{sec:parabolic}\label{ana:section}

The correction to the profile must satisfy the lower-order bound in
Lemma~\ref{geo:obstruction}.  To encode that bound together with its
spatial derivatives, put
\[
 \rho=\sqrt{-t},\qquad r=\sqrt{x},\qquad
 \sigma=(r^2+\rho)^{1/2},\qquad p_0=-2+2\delta,
 \qquad 0<\delta<1.
\]
At the concentrating point, a bound $|V|\le\eps\sigma^{p_0}$ gives
$|V(t,0)|\le\eps\rho^{-1+\delta}$, smaller than the profile's
$\rho^{-1}$ term.  On the sphere, $v_*\le r^2\le\sigma^2$, so it
also gives $|v_*V|\le\eps\sigma^{2\delta}$.  Since $\sigma$ is
bounded on $\Sigma\times[-T,0)$ for fixed $T$, a small weighted norm
will preserve positivity of the multiplier $1+v_*(F+V)$.
Outside the shrinking core, $r^2\ge\rho$,
the scale $\sigma$ is comparable to $r$ and the same weight records
spatial decay.  The particular value of $\delta$ will be chosen in
Theorem~\ref{flat:inverse}.

The norms below measure four spatial derivatives and one time derivative
in units $(\sigma,\sigma^4)$.  We also need estimates for an outer
forward problem with zero initial data, uniform as the elapsed time
tends to zero.  For this purpose we distinguish the full parabolic norm
from the norm that omits the first time derivative.

\subsection{Norms and scales}

On a product $D=U\times I\subset\R^d\times\R$, put
\[
 d_{\rm p}((z,t),(w,s))=|z-w|+|t-s|^{1/4}.
\]
Write
\[
 [g]_{\alpha;D}=\sup_{P\ne Q\in D}
       \frac{|g(P)-g(Q)|}{d_{\rm p}(P,Q)^\alpha},
 \quad |g|_{\alpha;D}=\|g\|_{\infty;D}+[g]_{\alpha;D}.
\]
The joint seminorm is equivalent, with constants depending only on
$\alpha$, to the sum of the spatial $\alpha$ and temporal $\alpha/4$
seminorms.  We use the following explicit version of the fourth-order
parabolic norm:
\begin{equation}\label{ana:local-norm}
 \mathcal N_\alpha(u;D)
   =\sum_{j=0}^4|D_z^ju|_{\alpha;D}+|u_t|_{\alpha;D},\qquad
 \mathcal N'_\alpha(u;D)=\sum_{j=0}^4|D_z^ju|_{\alpha;D}.
\end{equation}
Here and below norms of arrays mean sums of component norms.  Thus
$\mathcal N'_\alpha$ still includes temporal H\"older seminorms of all
four spatial derivatives.

\begin{definition}[Weighted local spaces]\label{ana:spaces}
Let $M$ be either $\R^m$ or the fixed sphere $\Sigma$, let $I$ be a time
interval, and let $s(P)>0$ be a scale on $M\times I$.  Fix a sufficiently
small number $a_0>0$.  In Euclidean coordinates, or normal coordinates
on $\Sigma$, centered at $P=(z_0,t_0)$, write $s_0=s(P)$ and set
\[
 \Psi_P(y,\theta)=(\exp_{z_0}(s_0y),t_0+s_0^4\theta),\qquad
 D_P=B_{a_0}\times
       \bigl((-a_0^4,a_0^4)\cap s_0^{-4}(I-t_0)\bigr).
\]
In the Euclidean case the exponential map means translation.  The
one-sided endpoints belonging to $I$ are included.  For a fixed weight
$p$, define
\begin{align}
 \|u\|_{X_{\alpha,p}(s)}
   &=\sup_P\mathcal N_\alpha(s_0^{-p}u\circ\Psi_P;D_P),\label{ana:X}\\
 \|u\|_{X'_{\alpha,p}(s)}
   &=\sup_P\mathcal N'_\alpha(s_0^{-p}u\circ\Psi_P;D_P),\label{ana:Xprime}\\
 \|f\|_{Y_{\alpha,p}(s)}
   &=\sup_P|s_0^{4-p}f\circ\Psi_P|_{\alpha;D_P}.
                                                        \label{ana:Y}
\end{align}
The derivatives in these formulas are derivatives of the displayed
pullbacks.  On the sphere $s$ is bounded above so that these charts,
and their fixed enlargements used below, lie within the injectivity
radius.  Restriction to radial functions on $\R^m$, or to invariant
functions on $\Sigma$, gives closed subspaces with these same norms.

Unless specified otherwise, $X,X',Y$ mean the spaces
\eqref{ana:X}--\eqref{ana:Y} with $s=\sigma$, $p=p_0$, and the
fixed exponent $0<\alpha<1$.  Subscripts $\alpha'$ indicate replacement
of $\alpha$ by $\alpha'$, with $\alpha<\alpha'<1$.
\end{definition}

These are local weighted norms; no H\"older quotient between distant
boxes is imposed.  Their completeness follows by convergence of the
functions and indicated derivatives on every compact subset, followed
by taking the supremum in the defining inequalities.  The derivative
relations pass to the limit distributionally.  In particular $X$ and
$X'$ are Banach spaces and $X\hookrightarrow X'$ has norm at most one.

The radius function $r$ is Lipschitz on both spaces.  The elementary
inequalities for Euclidean norms and fourth roots give
\[
 |\sigma(z,t)-\sigma(z_0,t_0)|
 \le |r(z)-r(z_0)|+|t-t_0|^{1/4}.
\]
Consequently $\sigma$ is comparable to its center value on every box
in Definition~\ref{ana:spaces}, including boxes truncated at $t=0$.
Changing $a_0$ to another fixed sufficiently small radius gives
equivalent norms: cover the enlarged box by finitely many smaller
comparable-scale boxes, and split a H\"older quotient whose points
are not in a common small box using the supremum norm.  The same
argument applies to $s_\beta=\max(\beta,r)$, whose lack of smoothness
at $r=\beta$ is harmless: the scale is frozen at each center, never
differentiated.  Products satisfy the usual local H\"older inequality
$|fg|_\alpha\le 2|f|_\alpha|g|_\alpha$, with their weights added.

\subsection{Interpolation without a lower bound for the time length}

We first isolate the interpolation used in localization.  For fixed
nested balls $B_1\Subset B_2$ and any interval $I$ of length at most
one, let
\[
 U=\|u\|_{\infty;B_2\times I},\qquad
 M=U+\|D_z^4u\|_{\infty;B_2\times I}
       +\|u_t\|_{\infty;B_2\times I}.
\]
Then, for $0\le j\le3$,
\begin{equation}\label{ana:interp-basic}
 |D_z^ju|_{\alpha;B_1\times I}
 \le C U^{1-(j+\alpha)/4}M^{(j+\alpha)/4}
 \le \eps M+C_\eps U.
\end{equation}
All constants are independent of $|I|$.  Spatial interpolation on
nested balls gives the spatial assertion, including the supremum.
For the temporal assertion apply that same interpolation to
$w=u(\cdot,t)-u(\cdot,s)$.  With $b=|t-s|$,
\[
 \|w\|_\infty\le\min(2U,bM),\qquad
 \|D^4w\|_\infty\le2M,
\]
and hence
\[
 \|D^jw\|_{\infty;B_1}
 \le C\min(2U,bM)^{1-j/4}M^{j/4}.
\]
Divide by $b^{\alpha/4}$ and split at $b=U/M$; both resulting
bounds give \eqref{ana:interp-basic}.  Restricting the available values
of $b$ can only decrease their supremum.  The spatial interpolation
inequalities themselves follow, for example, by subtracting a
Taylor polynomial in a ball of radius $q$ and minimizing in $q$;
fixed-radius remainders account for the $U$ term when the minimizer
exceeds the gap between the balls.  The same argument at order four
gives
\begin{equation}\label{ana:fourth-sup-interp}
 \|D^4u\|_{\infty;B_1\times I}
 \le\eps[D^4u]_{\alpha;B_2\times I}+C_\eps U.
\end{equation}
Here it suffices to use the spatial part of the seminorm on the
right.  Indeed, convolution on each spatial slice with a mollifier
of radius $q$ gives
$\|D^4u\|_{B_1}\le Cq^\alpha[D^4u]_{\alpha;B_2}
 +Cq^{-4}U$; choose $q$ below the spatial gap and so that
$Cq^\alpha\le\eps$.  Neither inequality uses a temporal seminorm
of $u_t$.

\subsection{Local fourth-order estimates}

Consider in a coordinate cylinder the real scalar operator
\begin{equation}\label{ana:operator}
 \mathcal L u=u_t+\sum_{|\nu|\le4}a_\nu(z,t)D^\nu u,
 \qquad
 \sum_{|\nu|=4}a_\nu(z,t)\xi^\nu\ge\lambda|\xi|^4.
\end{equation}
Assume $\lambda>0$ and $\sum_{|\nu|\le4}|a_\nu|_\alpha\le K$.
These assumptions are always imposed in the coordinates and units
of the estimate.  On a fixed smooth manifold they include the
coefficients arising from the coordinate expression of the operator.
The leading symbols used later are positive multiples of squared
positive definite quadratic forms, and satisfy \eqref{ana:operator}.

\begin{lemma}[Local Schauder estimates]\label{ana:schauder}
Under these assumptions the following estimates hold.
\begin{enumerate}[label=\textup{(\roman*)}]
\item If $Q_R=B_R\times(-R^4,0]$ and $\mathcal L u=f$ on $Q_2$, then
\begin{equation}\label{ana:interior-schauder}
 \mathcal N_\alpha(u;Q_1)
 \le C\bigl(|f|_{\alpha;Q_2}+\|u\|_{\infty;Q_2}\bigr).
\end{equation}
\item If $0<h\le1$, $\mathcal L u=f$ on $B_2\times[0,h]$, and
$u(\cdot,0)=0$, then
\begin{equation}\label{ana:initial-schauder}
 \mathcal N_\alpha(u;B_1\times[0,h])
 \le C\bigl(|f|_{\alpha;B_2\times[0,h]}
                   +\|u\|_{\infty;B_2\times[0,h]}\bigr).
\end{equation}
Here $C$ depends only on $d,\alpha,\lambda,K$ and the fixed spatial
gaps, and in \textup{(ii)} is independent of $h$.
\end{enumerate}
These are a priori estimates for classical solutions in the indicated
spaces; smooth approximation gives the same conclusions for limits
of such solutions.  The usual interior regularity theorem also
applies: with smooth coefficients, distributional solutions in
$L^2_{\rm loc}$ and H\"older sources are locally in the space
\eqref{ana:local-norm}, and are smooth when the source is smooth.
The estimates are from preceding times and require no terminal data.
\end{lemma}

\begin{proof}
We give the kernel and localization argument, particularly its
uniformity in \textup{(ii)}.  For a frozen homogeneous principal
operator $A_0$, with symbol $p(\xi)\ge\lambda|\xi|^4$, its heat kernel is
\[
 k_b(z)=(2\pi)^{-d}\int_{\R^d}e^{iz\cdot\xi-bp(\xi)}\,\dd\xi,
 \qquad b>0.
\]
Fourier differentiation and scaling give, uniformly for the allowed
symbols and every $N,j,\ell$,
\begin{equation}\label{ana:kernel}
 |D^j\partial_b^\ell k_b(z)|
 \le C_{N,j,\ell}b^{-(d+j)/4-\ell}
                  (1+|z|b^{-1/4})^{-N}.
\end{equation}
Moreover $\int k_b=1$ and $\int D^jk_b=0$ for $j>0$.  For
$u(t)=\int_0^tk_{t-s}*f(s)\,\dd s$ and $j=4$, subtract $f(s,z)$
inside the convolution.  Writing $F=\sup_s[f(s)]_{\alpha;\R^d}$,
\begin{equation}\label{ana:kernel-moment}
 \|D^4k_b*f(s)\|_\infty\le CF b^{-1+\alpha/4},\qquad
 \|\partial_bD^4k_b*f(s)\|_\infty
                    \le CF b^{-2+\alpha/4}.
\end{equation}
This proves $\|D^4u(t)\|_\infty\le CFt^{\alpha/4}$.
For a spatial increment of size $q$, split the lag integral at $q^4$.
The short-lag part is bounded by
$CF\int_0^{q^4}b^{-1+\alpha/4}\,\dd b$.
For the long-lag part use the mean value theorem and
$\|D^5k_b*f(s)\|_\infty\le CFb^{(\alpha-5)/4}$.
Its bound is $CFq\int_{q^4}^h b^{(\alpha-5)/4}\,\dd b$.
Both are at most $CFq^\alpha$, with an empty integral omitted.

For a temporal increment $e=t-s>0$, keep the source time unchanged:
\begin{align*}
 D^4u(t)-D^4u(s)
 ={}&\int_s^tD^4k_{t-r}*f(r)\,\dd r\\
 &+\int_0^s(D^4k_{t-r}-D^4k_{s-r})*f(r)\,\dd r.
\end{align*}
The first integral and the part of the second where $s-r<e$ are
bounded by $CF e^{\alpha/4}$.  On the remaining part,
\eqref{ana:kernel-moment} bounds the difference by
$CF e(s-r)^{-2+\alpha/4}$, and
\[
 e\int_e^s b^{-2+\alpha/4}\,\dd b\le C e^{\alpha/4}.
\]
This proves the full parabolic $\alpha$ bound for $D^4u$, down to
the initial face, without any requirement that $f(0)$ vanish.
The bounds for $u$ and lower derivatives follow from the kernel
integrals or \eqref{ana:interp-basic}; finally
$u_t=f-A_0u$ supplies the full parabolic norm of $u_t$.
We have proved the constant-coefficient zero-data estimate on every
slab of length at most one, with one constant.

Here are the localization details needed to preserve that constant.
We apply the fourth-derivative frozen estimate in the original local
units before localization; its constant is independent of the
localization radius.
In a spatial ball of radius $q$ and a time window of length at most
$q^4$, freeze the principal coefficients at one point.  Use a spatial
cutoff supported in that ball.  For windows meeting the initial
face use no lower-time cutoff, so the localized function still has
zero initial value.  For interior windows insert a time cutoff
which vanishes at the bottom and equals one on the smaller window;
its derivatives are bounded by fixed powers of $q^{-1}$.  The
constant-coefficient estimate then applies in both cases.  The
localized forcing consists of the original source, terms of spatial
order at most three, a time-cutoff multiple of $u$, and
$(a_\nu^0-a_\nu)D^\nu u$ for $|\nu|=4$.
The last terms obey
\begin{equation}\label{ana:freeze-bound}
 |(a_\nu^0-a_\nu)D^\nu u|_\alpha
 \le CKq^\alpha|D^4u|_\alpha
             +CK\|D^4u\|_\infty,
\end{equation}
with fixed cutoff factors understood.  First make $q$ small enough
to absorb the first coefficient.  Apply
\eqref{ana:fourth-sup-interp} to the second term.  For the lower-order
terms apply \eqref{ana:interp-basic}, and use the equation only in
the form
\begin{equation}\label{ana:ut-sup}
 \|u_t\|_\infty\le\|f\|_\infty
                   +K\sum_{j=0}^4\|D^ju\|_\infty.
\end{equation}
Thus no unknown temporal H\"older seminorm of $u_t$ is used to
absorb a temporal seminorm of $D^4u$.

For completeness, the resulting nested-domain absorption can be
expressed as follows.  Let $H(r)$ be the supremum and parabolic
$\alpha$ norm of $D^4u$ on the spatial ball of radius $r$, with the
given time interval in the initial case; in the interior case also
increase the preceding-time window with $r$.  For $1\le r<R\le3/2$,
localization, interpolation, and covering give, for any fixed
$\vartheta>0$,
\[
 H(r)\le\vartheta H(R)
       +C_\vartheta(R-r)^{-k}
                      (|f|_\alpha+\|u\|_\infty).
\]
Here $k$ is fixed; all losses are powers of spatial cutoff radii,
and none is a power of $h^{-1}$.  For H\"older pairs not sharing a
small localization window, the quotient is bounded by the
corresponding inverse power of that window's radius times
$\|D^4u\|_\infty$, which is handled by
\eqref{ana:fourth-sup-interp}.  Iterate along radii increasing
geometrically to $3/2$, choosing $\vartheta 2^k<1$.
The last term tends to zero since the solution norm is finite on
that compact cylinder.  This proves the estimate for $D^4u$.
Equations~\eqref{ana:interp-basic} and \eqref{ana:ut-sup} then give
the lower derivatives, and the equation in full H\"older norm gives
$|u_t|_\alpha$.  This proves \textup{(i)} and \textup{(ii)}.

For the distributional regularity assertion, extend the coefficients
and a spatially localized source from a compact coordinate cylinder
to a flat torus, combining the principal symbol with $|\xi|^4$
outside the coordinate patch.  This convex combination preserves
positivity.  The closed-manifold Cauchy theorem
\cite[Theorem~2.4]{Huang2026}, with order $r=4$ and zero data at
an earlier time, provides a particular solution
$w\in C^{4+\alpha,1+\alpha/4}$ near the smaller cylinder.
There $u-w$ is a homogeneous distributional solution.
To apply \cite[Theorem~5.1]{Hormander1961}, write the full symbol as
$p=i\tau+a_4(x,t,\xi)+a_{\le3}(x,t,\xi)$ and set
$R=1+|\tau|+|\xi|^4$.  On a compact coordinate set, ellipticity
and lower-order absorption give $|p|\ge cR$ at large frequencies.
Smoothness and polynomial degree then give
\[
 |\partial_{x,t}^{B}\partial_{\xi,\tau}^{A}p|
 \le C_{A,B}(1+|\xi|+|\tau|)^{-|A|/4}|p|.
\]
This verifies condition HE in
\cite[Equations~(3.1)' and~(3.2)]{Hormander1961}
with $d=1/4$ and $M_j=1$.  Bounded frequencies are covered by
$\widetilde p=(\sum_A|\partial_{\xi,\tau}^Ap|^2)^{1/2}
\ge|\partial_\tau p|=1$.
Thus $u-w$ is smooth, proving local H\"older regularity.
For smooth forcing the same theorem directly gives smoothness of $u$.

The fixed-cylinder
Schauder estimate, including an $L^2$ input for the solution, can
also be taken directly from \cite[Theorem~2.2\textup{(I)}]{DongZhang2015}.
Its estimate for the top spatial derivatives is parabolic;
the full temporal norm of $u_t$ follows here from the equation
and the assumed temporal H\"older regularity of the coefficients
and source.  This invocation concerns fixed interior cylinders;
the uniform initial-face assertion was proved above.
\end{proof}

\begin{lemma}[Interior estimate with an $L^2$ input]\label{ana:l2}
Let $U_1\Subset U_2$ be fixed bounded spatial domains and
$b_0<b_1<b$, with $b-b_1$ bounded above.  Suppose that the
coefficient assumptions of Lemma~\ref{ana:schauder} hold on
$U_2\times(b_0,b)$ and that $\mathcal Lu=f$ there.  Then
\begin{equation}\label{ana:l2-estimate}
 \mathcal N_\alpha(u;U_1\times[b_1,b))
 \le C\bigl(|f|_{\alpha;U_2\times(b_0,b)}
                  +\|u\|_{L^2(U_2\times(b_0,b))}\bigr).
\end{equation}
The constant depends on the spatial gap, $b_1-b_0$, the upper
bound for $b-b_1$, and the coefficient constants, but requires
no terminal trace.  The estimate is unchanged if the upper face
is included and the solution is defined there.  In particular it
holds on fixed nested annuli through an open upper time endpoint.
\end{lemma}

\begin{proof}
We recall an interpolation proof of the replacement of the supremum
term.  On three nested interior cylinders, with a gap of parabolic
size $a$, the estimate already proved bounds $\|D_zu\|_\infty$
and $\|u_t\|_\infty$ on the middle cylinder by
$Ca^{-k}(\|u\|_{\infty;\mathrm{outer}}+|f|_\alpha)$.
Average $u$ on a preceding cylinder of radius $q<a/4$ ending at
the point being estimated.  Its volume is a fixed multiple of
$q^{d+4}$, so
\[
 |u(z,t)|\le Cq^{-(d+4)/2}\|u\|_{L^2(\mathrm{outer})}
           +C(q+q^4)a^{-k}
                    (\|u\|_{\infty;\mathrm{outer}}+|f|_\alpha).
\]
Choose $q$ to be a sufficiently small fixed multiple of a suitable
power of $a$.  The last supremum coefficient can be made any
prescribed $\vartheta>0$, at the cost of another fixed inverse
power of $a$ on the other terms.  Iteration over geometrically
decreasing gaps, with $\vartheta$ smaller than the corresponding
geometric growth factor, bounds the inner supremum by
$C(\|u\|_{L^2}+|f|_\alpha)$.  Lemma~\ref{ana:schauder} on one
further nested cylinder proves \eqref{ana:l2-estimate}; finite
covering handles $U_1$ and the time interval.  All averaging
cylinders use preceding times.  First carry this out up to any
closed upper face $b'<b$ on which the local norms are finite,
then let $b'\uparrow b$.  The constants use only the fixed lower
time buffer.  Pairs of points approaching $b$ are included by
this limiting argument.
\end{proof}

\subsection{The outer forward problem}

The following formulation permits the localization scale to tend
to zero.  Fix a compact smooth manifold without boundary and a
time-independent positive Lipschitz scale $s(z)$.  Assume that on
balls of radius $4a_0s(z)$ the scale is comparable to its center
value, with uniform comparison constants and uniformly bounded
coordinate changes in scaled units.  We impose the coefficient
assumptions of \eqref{ana:operator} uniformly after scaling space
by $s(z)$ and time by $s(z)^4$; an original coefficient of spatial
order $j$ is thus multiplied by $s(z)^{4-j}$.  The ellipticity and
H\"older constants in these units are denoted again by $\lambda,K$.
These hypotheses hold for $s=s_\beta$ in the later application.

\begin{lemma}[Uniform zero-data forward estimate]\label{ana:short}
Fix $p\in\R$, $0<\alpha<1$, and the preceding scale and coefficient
constants.  There exist $\eta_0>0$ and $C<\infty$, independent of
$\inf s$, with the following property.  If
\[
 0<T\le\eta_0(\inf s)^4,\qquad
 \mathcal Lu=f\quad\hbox{on }M\times[0,T],\qquad u(\cdot,0)=0,
\]
then, with $\eta=T/(\inf s)^4$,
\begin{equation}\label{ana:short-bounds}
 \|u\|_{X_{\alpha,p}(s)}\le C\|f\|_{Y_{\alpha,p}(s)},\qquad
 \sup_{z,t}s(z)^{-p}|u(z,t)|
                         \le C\eta\|f\|_{Y_{\alpha,p}(s)}.
\end{equation}
For each fixed positive scale and smooth coefficients the zero-data
Cauchy problem has a unique solution with this regularity for
$f\in Y_{\alpha,p}(s)$.  The same estimates hold on every interval
$[0,T')$, $T'\le T$, on which the coefficients and source obey
the stated bounds.
\end{lemma}

\begin{proof}
Choose $\eta_0<a_0^4$.  Every defining box meets the initial
face, and its full time interval has length at most $\eta_0$
in center-scale units.  Apply \eqref{ana:initial-schauder} on a
fixed spatial enlargement of each such box, rescaling its fixed
spatial radius if necessary.  Scale comparability and the covering
observation after Definition~\ref{ana:spaces} give
\[
 \|u\|_{X_{\alpha,p}(s)}
 \le C\bigl(\|f\|_{Y_{\alpha,p}(s)}+S\bigr),\qquad
 S=\sup_{z,t}s(z)^{-p}|u(z,t)|.
\]
The pointwise time-derivative bound in this norm and zero initial
data imply
\[
 S\le\sup_z\frac{T}{s(z)^4}\|u\|_{X_{\alpha,p}(s)}
   \le\eta\|u\|_{X_{\alpha,p}(s)}.
\]
Decrease $\eta_0$ so that $C\eta_0\le1/2$, and absorb.  This
proves both estimates in \eqref{ana:short-bounds} with constants
independent of the number or size of the spatial boxes.

Fixed-scale existence and uniqueness follow from
\cite[Theorem~2.4]{Huang2026}, with the trivial real line bundle,
$r=4$, exponent $\alpha$, $L_t=-A(t)$, and zero initial data,
where $\mathcal L=\partial_t+A(t)$.
Its solution is $C^{4+\alpha,1+\alpha/4}$ on the closed time slab,
without requiring $f(\cdot,0)=0$; the equation gives
$u_t(\cdot,0)=f(\cdot,0)$.
Constants in this fixed-scale existence argument need not be
uniform in $\inf s$; uniformity of the asserted estimates has
already been established locally.  Truncation before an open
terminal face and passage to the limit prove the last statement.
\end{proof}

\begin{lemma}[Small spatial derivatives on a short interval]
\label{ana:interpolation}
Under the hypotheses of Lemma~\ref{ana:short}, the local scaled
parabolic $C^\alpha$ norms of the spatial derivatives of order
$0\le j\le3$, with weight $p-j$, satisfy
\begin{equation}\label{ana:lower-small}
 \sup_P\bigl|D_y^j(s_0^{-p}u\circ\Psi_P)\bigr|_{\alpha;D_P}
 \le C\eta^{(4-j-\alpha)/4}\|f\|_{Y_{\alpha,p}(s)}.
\end{equation}
If the coefficient and source bounds also hold at an exponent
$\alpha'\in(\alpha,1)$, then
\begin{equation}\label{ana:fourth-small}
 \|u\|_{X'_{\alpha,p}(s)}
 \le C\eta^{(\alpha'-\alpha)/4}
                    \|f\|_{Y_{\alpha',p}(s)}.
\end{equation}
Constants may depend on $\alpha'$ but remain independent of the
scale and elapsed time.
\end{lemma}

\begin{proof}
On a fixed enlarged scaled box, \eqref{ana:short-bounds} gives
$U\le C\eta\|f\|_Y$ and $M\le C\|f\|_Y$ in the notation of
\eqref{ana:interp-basic}.  That inequality proves
\eqref{ana:lower-small}.  In particular, all derivatives entering
a commutator of a fourth-order spatial operator with a spatial
cutoff have a common sufficient small factor
$\eta^{(1-\alpha)/4}$.

For the second claim apply Lemma~\ref{ana:short} at exponent
$\alpha'$, decreasing $\eta_0$ if necessary.  Every spatial
derivative of $u$ through order four has zero initial trace.
Indeed the solution and these derivatives are continuous to the
initial face, and the trace function $u(\cdot,0)$ is identically
zero.  For each $g=D_y^j(s_0^{-p}u\circ\Psi_P)$, $0\le j\le4$,
the temporal $\alpha'/4$ seminorm and the elapsed scaled time
therefore give
\[
 \|g\|_\infty\le C\eta^{\alpha'/4}
                            \|f\|_{Y_{\alpha',p}(s)},\qquad
 [g]_{\alpha'}\le C\|f\|_{Y_{\alpha',p}(s)}.
\]
On any metric space, splitting distances where the estimates
$|g(P)-g(Q)|\le2\|g\|_\infty$ and
$|g(P)-g(Q)|\le[g]_{\alpha'}d(P,Q)^{\alpha'}$ agree proves
\[
 [g]_\alpha\le
 C\|g\|_\infty^{1-\alpha/\alpha'}
                   [g]_{\alpha'}^{\alpha/\alpha'}.
\]
Apply this with the parabolic metric and take the supremum over
boxes.  It proves \eqref{ana:fourth-small}, including temporal
H\"older quotients of $D^4u$ with one point at the initial face.
\end{proof}

\begin{remark}\label{ana:no-time-smallness}
The exponent gap and omission of $u_t$ in
\eqref{ana:fourth-small} are necessary.  On a circle the exact
zero-data solution
\[
 u_N(z,t)=N^{-4-a}(1-e^{-N^4t})\cos(Nz),\qquad
 f_N(z)=N^{-a}\cos(Nz),\qquad T_N=N^{-4},
\]
of $u_t+\partial_z^4u=f_N$ has
$[\partial_z^4u_N]_{\alpha}$ of order $N^{\alpha-a}$ on
$[0,T_N]$.  This need not tend to zero when $a=\alpha$, whereas
$a=\alpha'$ gives precisely the power in
\eqref{ana:fourth-small}.  Moreover $u_t(0)=f_N$.
Thus the argument gives smallness of the full spatial norm $X'$,
while making no smallness assertion for the time derivative.
\end{remark}
